\section{Readout-Relative Lemmas and Structural Lifts}
\label{sec:readout-relative-lemmas}

The present framework distinguishes sharply between the structural kernel
of the theory and the visible form in which that kernel is read in a
given regime or frame. This distinction is not optional: the same
closure-stable object may appear as a probe/Dirac-type law, as an
effective QFT-type law in a phase-shifted observer regime, or as an
Einstein-type law after large-scale geometric readout. Consequently,
one must distinguish between \emph{kernel-level truths} and
\emph{readout-relative truths}.

The concrete instantiation of this distinction is already present in the
preceding spectral-projector block: the eigenvalue $\lambda_0$ is a
kernel-level invariant, while the measured value
$\lambda_0^{(\mathrm{matter})} = \gL^{-k}\cdot\lambda_0^{(\mathrm{photon})}$
(Remark~\ref{rem:frame-shift-spectral}) is readout-relative, depending
on the frame-insertion count $k$ of the observable in the matter frame
at $\zminus = +\bsym$.

\subsection*{Kernel level and readout level}

\begin{definition}[Kernel-level statement]
\label{def:kernel-level-statement}
A statement is called \emph{kernel-level} if it is invariant under the
admissible structural equivalences of the framework, including:
\begin{enumerate}[label=(\roman*)]
  \item sector-preserving unitary equivalence;
  \item compensator projection;
  \item observer-frame phase shift;
  \item admissible branch restriction.
\end{enumerate}
Such statements express properties of the underlying closure-stable
structural object itself rather than of one distinguished reading of it.
\end{definition}

\begin{definition}[Readout-relative lemma]
\label{def:readout-relative-lemma}
A statement $L^{(f)}$ is called a \emph{readout-relative lemma} if it
is proven only in a distinguished regime, frame, or readout $f$, and its
visible form depends on that readout. In particular, a readout-relative
lemma need not remain true as a frame-blind absolute statement.
\end{definition}

\begin{remark}
This is not a defect of the method. It is the correct formal response to
the fact that different regime-readings expose different visible shadows
of one common kernel. In the present framework this applies in particular
to the probe, effective QFT, and Einstein readouts.
The Symmetric matter--tachyon frame theorem (Theorem~\ref{thm:symmetric-frame})
is an example: $\beta_{\mathrm{matter}} = +\bsym$ is readout-relative
(it is the matter-frame reading), while the symmetry
$\beta_{\mathrm{tachyon}} = -\bsym$ and the common kernel $\bsym$ are
kernel-level invariants.
\end{remark}

\subsection*{Why this distinction is necessary}

The need for readout-relative lemmas is already built into the
architecture. The master system is closure-stable and unique only up to
the allowed sector-preserving, unitary, and frame transformations, while
the effective regimes differ by compensator decomposition, large-scale
suppression of closure fluctuations, and geometric readout. Thus one and
the same kernel may generate visibly different but structurally
compatible statements. In the present release this distinction should be read together with the OP transport theorem and the visible/structural mismatch separation tool: visible disagreement is secondary unless native closure-class preservation itself fails.

\begin{proposition}[Legitimacy of restricted lemmas]
\label{prop:legitimacy-restricted-lemmas}
A lemma proven in a distinguished readout $f$ may be used in the general
argument provided that:
\begin{enumerate}[label=(\roman*)]
  \item its frame or regime of validity is stated explicitly;
  \item its invariant kernel content is separated from its visible
        readout form;
  \item the passage from the readout-relative statement to the general
        claim is justified by a structural transport principle.
\end{enumerate}
\end{proposition}

\begin{proof}
The framework already distinguishes the common structural equation from
its regime projections (Theorem~\ref{thm:unity}, unitary equivalence via
Stone--von Neumann). Therefore a proof in one readout need not be
discarded merely because its visible form changes in another readout.
What must be transported is not the literal visible formula, but the
kernel-level content preserved under the admissible transport operations
of Lemmas~\ref{lem:group-transport} and~\ref{lem:hc-admissibility-lemma}.
\end{proof}

\subsection*{Structural lifts}

\begin{definition}[Structural lift]
\label{def:structural-lift}
Let $L^{(f)}$ be a readout-relative lemma valid in readout $f$.
A \emph{structural lift} of $L^{(f)}$ is a kernel-level statement
$L_{\mathrm{struct}}$ obtained by extracting from $L^{(f)}$ only the
invariant content that survives admissible unitary, sector, and frame
transport.
\end{definition}

\begin{remark}
A structural lift is therefore weaker than a literal universalization of
$L^{(f)}$, but stronger than leaving the lemma imprisoned in one special
case. The corollary Cor.~\ref{cor:SM-spectral-realization} is a concrete
example: statement~(iii) there lifts the frame-corrected eigenvalue
$\lambda_0^{(\mathrm{matter})}$ to the kernel-level invariant $\lambda_0$
by identifying $k=0$ for the singlet and $k\ge 1$ for the triplet.
\end{remark}

\begin{proposition}[Structural lift principle]
\label{prop:structural-lift-principle}
Suppose $L^{(f)}$ is proven in readout $f$, and suppose its proof
depends only on:
\begin{enumerate}[label=(\roman*)]
  \item the closure-stable master structure;
  \item invariants preserved under admissible transport;
  \item readout-specific visibility factors that can be isolated.
\end{enumerate}
Then there exists a structural lift $L_{\mathrm{struct}}$ that is valid
at kernel level, even if the visible statement $L^{(f)}$ itself is not
frame-blindly universal.
\end{proposition}

\begin{proof}
\noindent\textbf{Invariant content is transport-stable.}
By hypotheses~(i) and~(ii), the kernel-invariant content of the proof
of $L^{(f)}$ is preserved by the admissible transport operations.
Concretely: Lemma~\ref{lem:group-transport} guarantees that grouped-layer
transport carries admissible structure without creating new ontology, and
Lemma~\ref{lem:hc-admissibility-lemma} guarantees that compensator action
leaves the admissible sector invariant. Hence the kernel-invariant content
$P_{\mathrm{inv}}(\text{proof of }L^{(f)})$ is valid in every admissible
readout, not only in $f$.

\smallskip
\noindent\textbf{Readout-specific factors are isolatable.}
By hypothesis~(iii), the readout-dependent terms form a separable complement.
Discarding them while retaining $P_{\mathrm{inv}}(\text{proof})$ yields
$L_{\mathrm{struct}}$, independent of $f$.

\smallskip
\noindent\textbf{$L_{\mathrm{struct}}$ is kernel-level.}
Since $L_{\mathrm{struct}}$ depends only on the closure-stable structure
and transport-preserved invariants, it satisfies all four conditions of
Definition~\ref{def:kernel-level-statement}: invariance under
sector-preserving unitary equivalence, compensator projection,
observer-frame phase shift, and admissible branch restriction.
\end{proof}

\subsection*{The three-step proof logic}

This principle simplifies the global architecture considerably.
Instead of demanding that every intermediate statement already be valid
as a frame-blind absolute theorem, one may work in three steps:
\[
  \text{readout-relative lemma}
  \;\Longrightarrow\;
  \text{structural lift}
  \;\Longrightarrow\;
  \text{general kernel-level theorem}.
\]
This is the natural proof logic for a theory whose visible regimes are
different projections of one and the same structural object.

\subsection*{Shadow versus core: consequence for the compensator}

\begin{remark}[HC shadow versus HC kernel]
\label{rem:hc-shadow-vs-core}
This distinction is particularly important for the compensator.
In the minimal formulation, the visible compensator functional
$\mathcal{H}_C^{(f)}$ behaves as a closure-mismatch penalty in a given
readout and may vanish in an exact sector ($\mathcal{H}_C = 0$ in the
photon frame; cf.\ the closed master system). But under the
frame-dependent readout hypothesis developed above, such vanishing
concerns only the \emph{visible defect shadow} in a given regime, not
the deeper compensator-generating kernel
$C = \Pi^{(\mathcal{M})}_{\mathrm{spec}}$
(Theorem~\ref{thm:hc-spectral-projector}).

Hence a statement about $\mathcal{H}_C^{(f)}$ must not automatically be
read as a statement about the entire HC architecture. In particular, the
vanishing of the visible compensator field in the photon-frame readout
is consistent with the HC kernel remaining active as the
spectral-projector mechanism that selects the closure-stable sector ---
it is the shadow, not the core, that vanishes.

This is directly relevant to OP~5 ($\Lambda_{\mathrm{obs}}$): the
cosmological constant appears as a \emph{readout-relative} quantity ---
the frame offset $\bsym^2\langle f,f\rangle$ visible in the matter frame
--- while the photon-frame value $\Lambda_{\mathrm{fundamental}} = 0$
is the kernel-level statement. The large observed value is not a
fine-tuning problem but a failure to separate readout-relative visibility
from structural invariance.
\end{remark}

% ------------------------------------------------------------------
% Bidirectional Probe Lensing (v2.27.4 integration)
% ------------------------------------------------------------------

